\section{Polne mreže, negibne točke in izrek Knaster-Tarskega}

Naslednje vprašanje je, kakšni so zadostni pogoji za odsekoma linearno funkcijo, da ima negibno točko. Nedvoumno je, da to ne drži za vse odsekoma linearne preslikave. Preprost protiprimer je funkcija 
\[
f \colon \R \to \R,\quad f(x) = x + 1.
\]
%
Obenem znana izreka o negibnih točkah, kot sta Banachovo skrčitveno načelo in Brouwerjev izrek o negibni točki, tu ne pomagata, saj so odsekoma linearne funkcije v splošnem nezvezne. Odgovor na naše vprašanje bomo našli z uporabo teorije urejenosti.

Preden formuliramo osnovni izrek, navedemo osnovni definiciji. Vsebino tega razdelka povzemamo po \cite[1.\ poglavje]{zbMATH01576861}.

\begin{definicija}
    Naj bo \((E, \leq)\) delno urejena množica:
    \begin{itemize}
        \item \((E, \leq)\) je \emph{mreža}, če za poljubna elementa \(x, y \in E\), ima množica \(\set{x, y}\) največjo spodnjo mejo \(x \land y\) in najmanjšo zgornjo mejo \(x \lor y\). 
        \item \((E, \leq)\) je \emph{polna mreža}, če vsaka podmnožica \(X \subseteq E\) ima  največjo spodnjo mejo \(\bigwedge X\) in najmanjšo zgornjo mejo \(\bigvee X\). 
    \end{itemize}
\end{definicija}

\begin{definicija}
    Naj bosta \((E, \leq_E)\) in \((F, \leq_F)\) delno urejeni množici. Preslikava \(f \colon E \to F\) je \emph{monotona}, če
    \[
        \all{x, y \in E} x \leq_E y \lthen f(x) \leq_F f(y).
    \]
\end{definicija}

\begin{opomba}
  Beseda "`monotonost"' v tem delu ima enak pomen kot beseda "`naraš\-čajočost"'.
\end{opomba}

Sedaj lahko podamo osnovni izrek, ki bo dal odgovor na zastavljeno vprašanje.
\begin{izrek}[Knaster-Tarski]
  \label{izr:knaster-tarski}
  Naj bo \((E, \leq)\) polna mreža in \(f \colon E \to E\) monotona preslikava. Tedaj \(\bigwedge \Fix(f)\) in  \(\bigvee \Fix(f)\) sta elementi \(\Fix(f)\), kjer je 
  \[
    \Fix(f) = \setb{x \in E}{f(x) = x}.
  \]
\end{izrek}

Vidimo, da bo za monotono odsekoma linearno funkcijo 
\[
  f \colon E \subseteq \R \to E,
\]
kjer je \(E\) polna mreža, množica \(\Fix(f)\) neprazna in bo vsebovala minimum ter maksimum (torej najmanjšo in največjo negibno točko).  

Ustrezna izbira za množico \(E\) je na primer zaprt interval \([0,1]\), ki je očitno polna mreža, saj \(\R\) izpolnjuje Dedekindov aksiom.

Ko vemo, za kateri odsekoma linearni funkciji lahko izračunamo njene negibne točke, lahko problem formuliramo tudi v algoritmični obliki.